\documentclass[12pt, a4paper]{article}

\usepackage[utf8]{inputenc}
\usepackage[T1]{fontenc}
\usepackage[spanish]{babel}
\usepackage{geometry}
\usepackage{array}
\usepackage{booktabs}
\usepackage{tabularx}
\usepackage{multirow}
\usepackage{makecell}
\usepackage{enumitem}
\usepackage{fancyhdr}
\usepackage{amsmath, amssymb}
\usepackage{parskip}
\usepackage{titlesec}
\usepackage{rotating}
\usepackage{hyperref}

\geometry{top=2.5cm, bottom=2.5cm, left=2.8cm, right=2.5cm}
\hypersetup{colorlinks=false, hidelinks}

\titleformat{\section}{\large\bfseries}{\thesection.}{0.5em}{}[\titlerule]
\titleformat{\subsection}{\normalsize\bfseries}{}{0em}{}
\titlespacing*{\section}{0pt}{14pt}{6pt}
\titlespacing*{\subsection}{0pt}{10pt}{4pt}

\pagestyle{fancy}
\fancyhf{}
\fancyhead[L]{\small\textbf{Unidad Did\'actica --- Funciones Trigonom\'etricas}}
\fancyhead[R]{\small Exploraci\'on H\'aptica}
\fancyfoot[C]{\small\thepage}
\renewcommand{\headrulewidth}{0.4pt}
\setlength{\headheight}{15pt}

\renewcommand{\arraystretch}{1.3}

% Columna centrada verticalmente con ancho fijo
\newcolumntype{M}[1]{>{\centering\arraybackslash}m{#1}}
\newcolumntype{L}[1]{>{\raggedright\arraybackslash}m{#1}}

\begin{document}

% ============================================================
% PORTADA
% ============================================================
\begin{titlepage}
\centering

\vspace*{1.5cm}

{\large\bfseries Instituci\'on Educativa Alberto Lleras Camargo}\\[0.3em]
{\normalsize Colegio Alberto Lleras Camargo}

\vspace{3.5cm}

{\LARGE\bfseries UNIDAD DID\'ACTICA}\\[0.8em]
{\Large Funciones Trigonom\'etricas}\\[0.5em]
{\large\itshape Exploraci\'on h\'aptica de $\operatorname{sen}(x)$,
$\cos(x)$ y $\tan(x)$}

\vfill

\begin{tabular}{rl}
\textbf{Docentes en formaci\'on:} & Jhoan Alexander Enciso Romero \\[5pt]
                                   & Juan Sebasti\'an Pel\'aez Garc\'ia \\[5pt]
                                   & Daniel Mauricio Angulo Villabona \\[14pt]
\textbf{C\'odigos:}               & 141004710 \quad 141004620 \quad 141004616 \\[14pt]
\textbf{Jornada:}                 & \underline{\hspace{5cm}} \\[8pt]
\textbf{Grado:}                   & \underline{\hspace{5cm}} \\[8pt]
\textbf{Docente titular:}         & \underline{\hspace{5cm}} \\[8pt]
\textbf{Fecha:}                   & \underline{\hspace{5cm}} \\
\end{tabular}

\vspace{2.5cm}

\end{titlepage}

% ============================================================
\section{Justificaci\'on}
% ============================================================

El estudio de las funciones trigonom\'etricas exige que el estudiante construya
una imagen mental clara de curvas que oscilan, se repiten y, en el caso de la
tangente, presentan discontinuidades. Para estudiantes con baja o nula visi\'on,
los recursos visuales convencionales resultan inaccesibles sin una mediaci\'on
pedag\'ogica adecuada.

Esta unidad propone un canal de acceso alternativo: la \textbf{exploraci\'on
h\'aptica mediante l\'aminas en relieve}. El estudiante recorre las curvas con
los dedos, asociando cada textura y forma a un valor matem\'atico preciso. El
conocimiento construido es matem\'aticamente id\'entico al que obtiene un
estudiante vidente; \'unicamente el canal sensorial es diferente.

La unidad se fundamenta en el \textbf{Dise\~no Universal para el Aprendizaje
(DUA)}, el aprendizaje multisensorial (tacto + voz + movimiento), el andamiaje
progresivo y la promoci\'on de la autonom\'ia del estudiante.

% ============================================================
\section{Secuencia Metodol\'ogica (Plan)}
% ============================================================

\noindent\textbf{Duraci\'on total de la sesi\'on: 120 minutos (2 horas)}

\vspace{8pt}

% ------ TABLA PLAN -------
% Columnas: Etapa | Tiempo | Actividades | Objetivo | Conceptos | Estrategia | P.Investigacion
\setlength{\tabcolsep}{4pt}
\noindent
\begin{tabular}{|M{1.6cm}|M{1.3cm}|L{2.8cm}|L{2.4cm}|L{2.2cm}|L{2.0cm}|L{2.4cm}|}
\hline
\textbf{Etapas} &
\textbf{Tiempo} &
\textbf{Actividades} &
\textbf{Objetivo} &
\textbf{Conceptos involucrados} &
\textbf{Estrategias} &
\textbf{Pregunta de investigaci\'on} \\
\hline

% ---- FILA INICIAL 1 ----
\multirow{4}{1.6cm}{\centering\textbf{I\\N\\I\\C\\I\\A\\L}}
& 5 min
& Protocolarias: saludo, llamado a lista, pautas de convivencia.
& Crear clima de aula de confianza.
& ---
& Din\'amica de bienvenida.
& ?`C\'omo describirte a alguien, sin verla, una forma que s\'olo tus manos sienten? \\
\cline{2-7}

% ---- FILA INICIAL 2 ----
& 15 min
& Revisi\'on y desarrollo de tareas extraclase. Identificaci\'on de conductas de entrada.
& Detectar aprendizajes previos y dificultades pendientes.
& Funci\'on, \'angulo, valor de salida.
& Pregunta orientadora oral.
& ?`Qu\'e errores conceptuales persisten de la sesi\'on anterior? \\
\hline

% ---- FILA DESARROLLO 1 ----
\multirow{8}{1.6cm}{\centering\textbf{D\\E\\S\\A\\R\\R\\O\\L\\L\\O}}
& 10 min
& Ideas o conceptos previos. (Algunos identificados en conductas de entrada.)
& Activar saberes sobre oscilaci\'on y repetici\'on.
& Oscilaci\'on, periodicidad (intuitivo).
& Lluvia de ideas.
& ?`Con qu\'e claridad llegan al t\'ermino ``periodicidad''? \\
\cline{2-7}

% ---- FILA DESARROLLO 2 ----
& 15 min
& \textbf{Motivaci\'on:} introducci\'on verbal --- ?`Qu\'e vamos a sentir? Met\'afora de la colina y la ola.
& Motivar y contextualizar la exploraci\'on h\'aptica.
& Funci\'on, entrada, salida, oscilaci\'on.
& Relato guiado y met\'afora sensorial.
& ?`Relaciona el estudiante la met\'afora con el concepto matem\'atico? \\
\cline{2-7}

% ---- FILA DESARROLLO 3 ----
& 20 min
& \textbf{Exploraci\'on del seno.} Recorrido de la l\'amina $\operatorname{sen}(x)$. Conceptualizaci\'on y teorizaci\'on.
& Que el estudiante identifique m\'aximo, m\'inimo, per\'iodo y rango del seno.
& $\operatorname{sen}(x)$, per\'iodo, amplitud, rango.
& Exploraci\'on h\'aptica + verbalizaci\'on.
& ?`Asocia correctamente cada punto palpado con su valor num\'erico? \\
\cline{2-7}

% ---- FILA DESARROLLO 4 ----
& 20 min
& \textbf{Comparaci\'on seno/coseno.} Ambas l\'aminas simult\'aneas. Desfase de fase.
& Que el estudiante perciba y describa el desfase entre las dos funciones.
& $\cos(x)$, desfase, identidad trigonom\'etrica.
& Exploraci\'on h\'aptica comparativa.
& ?`Comprende que el coseno es el seno desplazado $\pi/2$? \\
\cline{2-7}

% ---- FILA DESARROLLO 5 ----
& 20 min
& \textbf{Exploraci\'on de la tangente.} Rampas, as\'intotas, discontinuidades.  Ampliaci\'on.
& Que el estudiante distinga una funci\'on no acotada con discontinuidades.
& $\tan(x)$, as\'intota, discontinuidad, rango infinito.
& Exploraci\'on h\'aptica + preparaci\'on verbal previa.
& ?`Comprende que la discontinuidad es inevitable y no un defecto de la l\'amina? \\
\hline

% ---- FILA CIERRE ----
\textbf{CIERRE}
& 20 min
& S\'intesis y conclusi\'on. Preguntas guiadas con la l\'amina. Tarea extraclase.
& Consolidar aprendizajes y asignar trabajo aut\'onomo.
& Todos los conceptos trabajados.
& Preguntas con respuesta h\'aptica.
& ?`Puede el estudiante relacionar el valor h\'aptico con el valor num\'erico de forma aut\'onoma? \\
\hline

\end{tabular}

% ============================================================
\section{Desarrollo del Plan}
% ============================================================

\textit{Todo lo siguiente est\'a escrito y resuelto seg\'un la estructura de la
unidad did\'actica.}

% -----------------------------------------------------------
\subsection{1. Actividades protocolarias \normalfont(5 min)}

\textbf{Estrategia:} din\'amica grupal de bienvenida.

El docente saluda usando el nombre de cada estudiante al pasar lista y establece
un ambiente de confianza. Se acuerdan las pautas de convivencia: hablar por
turnos, describir en voz alta lo que se palpa, no interrumpir la exploraci\'on
del compa\~nero.

\textbf{Recursos:} planilla de asistencia; l\'aminas h\'apticas guardadas a\'un
en la carpeta.

\textbf{Pregunta de investigaci\'on (actitudinal):} \textit{``?`C\'omo le
describir\'ias a alguien, sin verla, una forma que s\'olo tus manos pueden
sentir?''}

% -----------------------------------------------------------
\subsection{2. Revisi\'on de tarea extraclase \normalfont(15 min)}

\textbf{Proceso:} si la sesi\'on anterior dej\'o tarea, el docente solicita a
2 o 3 estudiantes compartir sus respuestas oralmente. Se identifican:

\begin{itemize}[itemsep=2pt]
  \item Dificultades en la comprensi\'on de ``entrada'' y ``salida'' de una
        funci\'on.
  \item Confusiones entre valor m\'aximo y m\'inimo.
  \item Uso impreciso de ``repetici\'on'' frente a ``periodicidad''.
\end{itemize}

\textbf{Pregunta de investigaci\'on:} \textit{``?`Qu\'e errores conceptuales o
procedimentales persisten desde la sesi\'on anterior y requieren refuerzo
hoy?''}

% -----------------------------------------------------------
\subsection{3. Ideas previas \normalfont(10 min)}

\textbf{Estrategia:} lluvia de ideas guiada.

El docente plantea oralmente tres preguntas sin corregir las respuestas en este
momento:

\begin{enumerate}[itemsep=3pt]
  \item ?`Qu\'e significa que algo sea \textit{peri\'odico}?
  \item ?`Conocen alg\'un fen\'omeno que suba y baje de forma regular?
  \item ?`Qu\'e esperan ``sentir'' en una l\'amina que representa una funci\'on?
\end{enumerate}

\textbf{Ideas previas esperadas:} onda, ritmo, ciclo, repetici\'on, m\'aximo,
m\'inimo.

\textbf{Pregunta de investigaci\'on:} \textit{``?`Con qu\'e claridad conceptual
llegan los estudiantes al t\'ermino `periodicidad'?''}

% -----------------------------------------------------------
\subsection{4. Presentaci\'on del tema}

\subsubsection*{Motivaci\'on --- ?`Qu\'e vamos a sentir? \normalfont(15 min)}

Antes de tocar las l\'aminas, el docente establece el marco con lenguaje
concreto:

\begin{itemize}[itemsep=3pt]
  \item Una \textbf{funci\'on} es una relaci\'on entre una entrada (el \'angulo
        $x$) y una salida (el valor de la funci\'on).
  \item Las funciones trigonom\'etricas \textbf{oscilan}: suben y bajan de forma
        repetida e indefinida.
  \item La l\'amina representa esa oscilaci\'on como \textbf{relieve f\'isico}:
        alto $=$ positivo; nivel del papel $=$ cero; hundido $=$ negativo.
\end{itemize}

\textbf{Met\'afora de inicio:}

\begin{quote}
\itshape ``Imaginen que caminan por una colina. Cuando ascienden, el valor de
la funci\'on crece. En la cima llegan al m\'aximo. Cuando descienden vuelven
a cero, y si siguen bajando entran al valle: el m\'inimo. Luego la colina
vuelve a comenzar. Eso es lo que van a recorrer con los dedos.''
\end{quote}

\textbf{Pregunta de investigaci\'on:} \textit{``?`Relaciona el estudiante la
met\'afora de la colina con la curva matem\'atica?''}

% - - - - - - - - - - - - - - - - - - - - - - - - - - - -
\subsubsection*{Conceptualizaci\'on y Teorizaci\'on ---
Exploraci\'on del Seno \normalfont(20 min)}

El estudiante recibe la l\'amina de $f(x)=\operatorname{sen}(x)$ y la recorre
con el dedo \'indice de izquierda a derecha. El docente verbaliza
simult\'aneamente:

\vspace{4pt}
\begin{tabular}{|p{6.5cm}|p{6.5cm}|}
\hline
\textbf{Lo que palpa} & \textbf{Lo que significa} \\
\hline
Parte desde el nivel medio        & $\operatorname{sen}(0)=0$ \\
\hline
Sube suavemente hasta la cresta   & $\operatorname{sen}(\pi/2)=1$ --- m\'aximo \\
\hline
Baja de vuelta al nivel medio     & $\operatorname{sen}(\pi)=0$ \\
\hline
Sigue bajando hasta el valle      & $\operatorname{sen}(3\pi/2)=-1$ --- m\'inimo \\
\hline
Regresa al nivel medio            & $\operatorname{sen}(2\pi)=0$ --- ciclo completo \\
\hline
La forma se repite                & La funci\'on es peri\'odica, $T=2\pi$ \\
\hline
\end{tabular}

\vspace{6pt}
\textbf{Teorizaci\'on:}
\[
  \operatorname{sen}(x):\quad
  \text{per\'iodo}=2\pi, \qquad
  \text{amplitud}=1, \qquad
  \text{rango}=[-1,\,1]
\]

\textbf{Aplicaci\'on:} el estudiante localiza sobre la l\'amina los puntos
$x=0$, $\pi/2$, $\pi$, $3\pi/2$ y $2\pi$, y dicta en voz alta el valor
correspondiente.

\textbf{Pregunta de investigaci\'on:} \textit{``?`Asocia correctamente cada
punto palpado con su valor num\'erico?''}

% - - - - - - - - - - - - - - - - - - - - - - - - - - - -
\subsubsection*{Conceptualizaci\'on --- Comparaci\'on seno / coseno
\normalfont(20 min)}

Con la l\'amina de $f(x)=\cos(x)$:

\begin{itemize}[itemsep=3pt]
  \item El estudiante nota que la curva \textbf{arranca desde la cresta}:
        $\cos(0)=1$, no desde el nivel medio.
  \item Se colocan ambas l\'aminas una junto a la otra.
  \item Con una mano en cada l\'amina, el estudiante palpa el \textbf{desfase}:
        donde el seno tiene su cresta ($x=\pi/2$), el coseno ya est\'a bajando.
\end{itemize}

\begin{quote}
\itshape ``Son la misma ola, pero una va adelantada respecto a la otra. Ese
adelanto equivale exactamente a un cuarto de per\'iodo: $\pi/2$ radianes.''
\end{quote}

\textbf{Teorizaci\'on:}
\[
  \cos(x)=\operatorname{sen}\!\Bigl(x+\tfrac{\pi}{2}\Bigr)
  \qquad\text{(desfase de }\pi/2\text{ radianes)}
\]

\textbf{Pregunta de investigaci\'on:} \textit{``?`Comprende el estudiante que
el coseno es el seno desplazado $\pi/2$ unidades hacia la izquierda?''}

% - - - - - - - - - - - - - - - - - - - - - - - - - - - -
\subsubsection*{Ampliaci\'on --- Exploraci\'on de la Tangente
\normalfont(20 min)}

La tangente requiere preparaci\'on verbal \textbf{antes} de tocar la l\'amina:

\begin{itemize}[itemsep=3pt]
  \item \textbf{No est\'a acotada:} no tiene m\'aximo ni m\'inimo finitos.
  \item \textbf{Tiene as\'intotas verticales} en $x=\dfrac{\pi}{2}+k\pi$
        ($k\in\mathbb{Z}$): la funci\'on se ``rompe'' y salta de
        $-\infty$ a $+\infty$.
\end{itemize}

Al palpar la l\'amina el estudiante siente:

\begin{enumerate}[itemsep=3pt]
  \item Rampas que suben muy abruptamente y se interrumpen.
  \item Un vac\'io o corte f\'isico donde est\'a la as\'intota.
  \item La curva aparece de nuevo desde abajo y vuelve a subir.
\end{enumerate}

\begin{quote}
\itshape ``Siente c\'omo sube sin parar hasta que desaparece\ldots{} y luego
vuelve a empezar desde muy abajo. No hay cresta: no hay techo.''
\end{quote}

\textbf{Teorizaci\'on:}
\[
  \tan(x)=\frac{\operatorname{sen}(x)}{\cos(x)},\qquad
  \text{no definida en }x=\frac{\pi}{2}+k\pi,\quad k\in\mathbb{Z}
\]
\[
  \text{per\'iodo}=\pi,\qquad\text{rango}=(-\infty,\,+\infty)
\]

\textbf{Pregunta de investigaci\'on:} \textit{``?`Comprende el estudiante que
la discontinuidad de $\tan(x)$ es matem\'aticamente inevitable y no un error
de la l\'amina?''}

% -----------------------------------------------------------
\subsection{Evaluaci\'on}

\textbf{Formativa --- preguntas guiadas con respuesta h\'aptica:}

El docente plantea las preguntas y el estudiante responde explorando con los
dedos sobre la l\'amina correspondiente:

\begin{enumerate}[itemsep=4pt]
  \item ?`Cu\'al funci\'on empieza en cero?
  \item ?`Cu\'al empieza en su valor m\'as alto?
  \item ?`Cu\'antos valles hay en un ciclo completo del seno?
  \item ?`En qu\'e punto se ubica la primera as\'intota de la tangente?
  \item ?`Cu\'al es la diferencia de fase entre el seno y el coseno?
  \item ?`Qu\'e funci\'on no tiene amplitud definida? ?`Por qu\'e?
\end{enumerate}

\textbf{Pregunta de investigaci\'on (actitudinal):} \textit{``?`Muestra el
estudiante seguridad al explorar la l\'amina de forma aut\'onoma, sin esperar
validaci\'on constante del docente?''}

\vspace{6pt}
\textbf{Sumativa --- taller individual} (entrega pr\'oxima sesi\'on):

\begin{enumerate}[itemsep=4pt]
  \item Con la l\'amina del seno, identifica y registra el valor de
        $\operatorname{sen}(x)$ para
        $x\in\{0,\;\pi/2,\;\pi,\;3\pi/2,\;2\pi\}$.
  \item ?`En qu\'e valores de $x$ son iguales $\operatorname{sen}(x)$ y
        $\cos(x)$? (Busca los puntos de cruce entre ambas l\'aminas.)
  \item Describe con tus palabras la diferencia entre una funci\'on peri\'odica
        acotada y una no acotada.
  \item \textbf{Profundizaci\'on:} investiga qu\'e relaci\'on tiene $\tan(x)$
        con la pendiente de una recta.
\end{enumerate}

\textbf{Pregunta de investigaci\'on (dominio):} \textit{``?`Puede el estudiante
relacionar el valor h\'aptico de la l\'amina con el valor num\'erico de la
funci\'on de forma aut\'onoma?''}

% ============================================================
\section{Actividad de Cierre \normalfont--- S\'intesis y conclusi\'on
\normalfont(20 min)}
% ============================================================

El docente retoma las ideas previas del inicio y las contrasta con lo aprendido:

\begin{itemize}[itemsep=3pt]
  \item ?`Cambi\'o su imagen mental de una ``ola''?
  \item ?`En qu\'e se parecen y en qu\'e se diferencian las tres funciones?
\end{itemize}

\textbf{Reflexi\'on escrita u oral:} cada estudiante completa (en tarjeta,
Braille o verbalmente):

\begin{quote}
\itshape ``Hoy aprend\'i que \underline{\hspace{4.5cm}}. Lo que m\'as me
sorprendi\'o fue \underline{\hspace{4cm}}. Todav\'ia tengo dudas sobre
\underline{\hspace{3.5cm}}.''
\end{quote}

\textbf{Tarea extraclase:} completar el taller sumativo de la secci\'on
anterior (entrega pr\'oxima sesi\'on). Para profundizaci\'on: lectura sobre la
aplicaci\'on de las funciones trigonom\'etricas en el an\'alisis de se\~nales
sonoras (nociones de series de Fourier).

\vspace{10pt}
\begin{tabular}{|p{14.5cm}|}
\hline
\vspace{2pt}
\textbf{Recursos y materiales}\\[4pt]
\begin{itemize}[itemsep=2pt, topsep=0pt]
  \item L\'aminas h\'apticas en relieve de $\operatorname{sen}(x)$,
        $\cos(x)$ y $\tan(x)$ (una por estudiante o en parejas).
  \item Planilla de asistencia (tinta y/o Braille seg\'un aplique).
  \item Tarjetas de reflexi\'on (cart\'on grueso).
  \item Regla t\'actil o cord\'on para identificar el eje horizontal.
  \item Material del taller en el formato accesible de cada estudiante
        (tinta, Braille o archivo digital con lector de pantalla).
\end{itemize}
\vspace{2pt}\\
\hline
\end{tabular}

\vspace{1.5cm}
\noindent\rule{\linewidth}{0.4pt}
\begin{center}
\small Elaborado con base en la \textit{Sesi\'on Did\'actica: Funciones
Trigonom\'etricas para Personas con Baja o Nula Visi\'on} y la
\textit{Estructura de Unidad Did\'actica} institucional.
\end{center}

\end{document}
